\section{Introducción: del furor por el concepto de vuelta a la productividad}

Esta sección establece primero la manera de leer el episodio completo. Wang He (王鹤) es a la vez profesor asistente de la Universidad de Pekín y fundador y CTO de Yinhe Tongyong (银河通用); su relato abarca al mismo tiempo la historia académica, las rutas tecnológicas, las decisiones de emprendimiento y la ética de la industria. Más que discutir solo «si los robots llegarán pronto a los hogares», las preguntas más importantes de este episodio son: por qué la inteligencia corporizada pasó de ser una línea marginal de la visión por computadora a convertirse en la favorita del capital, por qué el lenguaje no es la única puerta de entrada a la inteligencia, por qué los datos para VLM/VLA son más escasos que los datos para LLM y por qué, en última instancia, la industria debe demostrar su valor con productividad y no con demos.

Este episodio complementa la entrevista con Xie Chen (谢晨) del EP109. El EP109 explica «cómo se fabrican los datos» desde la infraestructura de simulación y de datos sintéticos; el EP106 explica, desde las escuelas académicas, el ciclo cerrado visión-acción y la lógica de implantación de una empresa, «por qué se necesita este tipo de datos y qué narrativas comerciales dañan a la industria». En conjunto, los dos episodios permiten articular la inteligencia corporizada como una cadena completa que va del concepto en los artículos académicos a las rutas de modelos, la recipe de datos, las restricciones de hardware y el ROI en escenarios reales.

\begin{importantbox}{Tesis central del episodio}
La inteligencia corporizada no es algo tan simple como «meter un modelo grande en un robot». En primer lugar, exige que la percepción, la acción y la retroalimentación del entorno formen un ciclo cerrado; en segundo lugar, exige que el hardware, los datos, los modelos y los escenarios avancen en espiral ascendente; por último, debe demostrar su valor con productividad escalable, y no con montos de financiamiento, demos de teleoperación o narrativas vagas sobre robots de propósito general.
\end{importantbox}

\begin{knowledgebox}{Criterio visual}
Este video es una entrevista con encuadre fijo, sin slides, pizarrón ni demostraciones de producto. El texto usa la portada solo para identificar la fuente; todas las imágenes del texto son diagramas conceptuales de elaboración propia, que sirven para explicar el lenguaje y la inteligencia, la genealogía académica de la inteligencia corporizada, el Perception-Action Loop, la recipe de datos, la productividad como producto y los excesos del capital.
\end{knowledgebox}

\subsection{Resumen de la sección}

El hilo conductor del EP106 es que la inteligencia corporizada pasó del margen académico de la visión al centro de la industria, pero el verdadero criterio no es qué tanto entusiasmo despierta la historia, sino si se logra poner en marcha al mismo tiempo el ciclo cerrado percepción-acción, el volante de datos y la productividad comercial.
